% %!TEX root = ../../main.tex

% \section{Notation and Conventions}

% Throughout, we use the following notation:
% \begin{itemize}
% \item $\norm{\cdot}$ denotes the operator (spectral) norm of matrices
% \item $\norm{\cdot}_F$ denotes the Frobenius norm: $\norm{M}_F^2 = \sum_{ij} |M_{ij}|^2$
% \item $N$ denotes the matrix dimension; when the Kronecker decomposition $N = qp$ is active (Section~\ref{sec:kronecker-skip}), $q$ is the twist dimension and $p$ is the fiber dimension
% \item $I$ or $I_N$ denotes the $N \times N$ identity matrix
% \item $\C^+ = \{z \in \C : \Im(z) > 0\}$ is the upper half-plane
% \item For $z \in \C$, we write $z = x + iy$ with $x = \Re(z)$ and $y = \Im(z)$
% \item $A_q \in \R^{q \times q}$: the \emph{twist matrix} in Hyper-Connections (Section~\ref{sec:kronecker-skip})
% \item $\mcB = M_q(\C)$: the base algebra for operator-valued free probability (Section~\ref{sec:ov-free-prob})
% \item $E_\mcB = \mathrm{id}_q \otimes \tr_p$: the $\mcB$-valued conditional expectation (normalized partial trace over $p$-fibers)
% \item $G^{(\mcB)}(z)$: the $\mcB$-valued Cauchy transform; $G(z) = \frac{1}{q}\Tr_q(G^{(\mcB)}(z))$ is its scalar projection
% \item $S^{(\mcB)}(W)$: the $\mcB$-valued S-transform (Section~\ref{sec:ov-s-transform})
% \item $A_h(M) = A_q + \sigma^2 M$: the \emph{dressed matrix} in the Dyson equation
% \end{itemize}

% \begin{definition}[Block Trace Operator]\label{def:block-trace}
For a $2N \times 2N$ matrix $M$ partitioned as
\[
M = \begin{pmatrix} M_{11} & M_{12} \\ M_{21} & M_{22} \end{pmatrix},
\quad M_{ij} \in \R^{N \times N},
\]
the \emph{block trace} is the $2 \times 2$ matrix of normalized traces:
\[
\bTr(M) = \begin{pmatrix} \frac{1}{N}\Tr(M_{11}) & \frac{1}{N}\Tr(M_{12}) \\ \frac{1}{N}\Tr(M_{21}) & \frac{1}{N}\Tr(M_{22}) \end{pmatrix}.
\]
\end{definition}


% \begin{definition}[Isotropy]\label{def:isotropy}
A random matrix $X \in \R^{N \times N}$ is \emph{isotropic} with variance parameter $\sigma^2$ if $\E[X] = 0$ and
\[
\E[XBX^\top] = \frac{\sigma^2}{N}\Tr(B) I \quad \text{for all deterministic } B \in \R^{N \times N}.
\]
\end{definition}


% \begin{example}[Isotropic ensembles]
% The following random matrix ensembles are isotropic:
% \begin{enumerate}
% \item \textbf{Gaussian:} $X_{ij} \sim \mathcal{N}(0, \sigma^2/N)$ i.i.d.
% \item \textbf{Orthogonally invariant:} $X = U \Lambda V^\top$ where $U, V$ are Haar-distributed on $O(N)$ and $\Lambda$ is diagonal with appropriate variance
% \item \textbf{Scaled Wigner:} $X = (X + X^\top)/(2\sqrt{N})$ with $X_{ij} \sim \mathcal{N}(0, \sigma^2)$ (symmetric isotropic ensemble)
% \end{enumerate}
% \end{example}
